\section{Local operators on the invariant boundary}%
\label{sec:peps_regular_boundary_local_algebra}

The invariant boundary of \cite[Theorem~6.9]{Schuch2010PEPS},
local source lines 2043--2076, imposes a constraint on which
operators on the original boundary legs can act on its support.
The following auxiliary classification concerns that constraint.
It gives no estimate on the spatial decay of a boundary Hamiltonian.

\begin{theorem}[The local algebra compatible with invariant support]%
    \label{thm:peps_regular_boundary_local_algebra}
    \lean{TNLean.PEPS.regularBoundaryLocalMap,
        TNLean.PEPS.regularBoundaryLocalMap_preserves_invariants_iff,
        TNLean.PEPS.regularLeg_commute_iff_entries,
        TNLean.PEPS.regularBoundaryLocalMap_preserves_invariants_iff_commute}
    \leanok
    \uses{thm:peps_regular_boundary_invariants}
    Let $I,J$ be disjoint finite sets of regular boundary legs,
    with $J\ne\varnothing$. Let $\mathcal H_I=\mathbb C^{G^I}$
    and $\mathcal H_J=\mathbb C^{G^J}$, and let $L_I(g)$ and
    $L_J(g)$ be the simultaneous left translations.
    For every linear operator $A$ on $\mathcal H_I$, the following
    conditions are equivalent:
    \begin{enumerate}
        \item $A\otimes I$ preserves the invariant subspace of
              $L_I\otimes L_J$;
        \item $A L_I(g)=L_I(g)A$ for every $g\in G$;
        \item $A(ga,gc)=A(a,c)$ for every $g\in G$ and $a,c\in G^I$.
    \end{enumerate}
    Self-adjointness of $A$ is not required.
    These local algebras concern the regular invariant boundary; see
    \cite{gap:rmp_peps_quantum_double_g_isometry}.
\end{theorem}

\begin{proof}\leanok
    For $c\in G^I$, the vector
    \[
        v_c=\sum_{g\in G}|gc\rangle\otimes|g,\ldots,g\rangle
    \]
    is globally invariant. The coefficient of
    $(A\otimes I)v_c$ at $(a,(g,\ldots,g))$ is $A(a,gc)$:
    the nonempty complementary set distinguishes the common
    translation $g$. If this image is invariant, comparison of
    its coefficients at $(a,(1,\ldots,1))$ and
    $(ga,(g,\ldots,g))$ gives $A(ga,gc)=A(a,c)$.
    Conversely, this matrix identity and a change of summation
    variable show that $A\otimes I$ preserves every invariant
    vector. Comparing the two compositions on coordinate vectors
    identifies the same matrix identity with commutation.
\end{proof}

\begin{corollary}[A local matrix-unit obstruction]%
    \label{cor:peps_regular_boundary_local_matrix_unit}
    \lean{TNLean.PEPS.not_regularBoundaryLocalMap_single_preserves_invariants}
    \leanok
    \uses{thm:peps_regular_boundary_local_algebra}
    Suppose $G$ is nontrivial and both $I,J$ are nonempty.
    The projector $|1,\ldots,1\rangle\langle1,\ldots,1|$
    on the selected legs, extended by the complementary identity,
    does not preserve the global invariant space.
    Thus the full matrix algebra on a nonempty proper subset
    of the original boundary legs does not restrict to that space.
\end{corollary}

\begin{proof}\leanok
    Choose $g\ne1$. The projector has diagonal entry one at the
    identity configuration and zero at its common translate by $g$.
    Apply the preceding classification.
\end{proof}
